%==============================================================================
\section{Materials and Methods}
\label{sec:methods}

\subsection{Dataset and splits}
\label{sec:methods-data}

We use the CAMUS dataset \citep{leclerc2019}, comprising 500 patients
imaged in apical-2-chamber (2CH) and apical-4-chamber (4CH) views, with
expert annotations at end-diastole (ED) and end-systole (ES). Each frame
is labelled with four classes: background, LV-endocardium, LV-epicardium,
and left atrium. Per-patient pixel spacing is recorded in the NIfTI
header and ranges between approximately $0.30$ and $0.42$~mm.

We adopt the official CAMUS split: $400$ patients for training,
$50$ for validation, and $50$ for testing, with no overlap. After
ED+ES frame extraction this yields approximately $1{,}600$ training,
$200$ validation, and $200$ test images. Half-cycle sequences are
optionally used for augmentation during training; results are reported
on the canonical ED+ES test set only.

\subsection{Architectures}
\label{sec:methods-archs}

We benchmark nine representative architectures spanning three paradigms:

\paragraph{Pure CNN (seven architectures).}
\begin{itemize}[leftmargin=*,nosep]
\item \textbf{UNet-V1} \citep{ronneberger2015}: canonical encoder--decoder
      with skip connections and trainable transposed-convolution
      upsampling.
\item \textbf{UNet-V2} (custom): residual blocks, squeeze-and-excitation
      \citep{hu2018}, attention-gated skips \citep{oktay2018}, and
      optional deep supervision.
\item \textbf{UNet-ResNet}: U-Net with a ResNet-34 encoder
      \citep{he2016} pretrained on ImageNet.
\item \textbf{DeepLabV3+} \citep{chen2018encoder}: ResNet-50 encoder,
      atrous spatial pyramid pooling bottleneck, lightweight decoder.
\item \textbf{nnU-Net} \citep{isensee2021}: auto-configured U-Net
      variant, trained from scratch.
\item \textbf{DenseContextU-Net} (custom): dense blocks
      \citep{huang2017} with a global-context module in the bottleneck.
\item \textbf{FPN-UNet}: U-Net decoder with a Feature Pyramid
      \citep{lin2017} neck over a ResNet-50 encoder.
\end{itemize}

\paragraph{Hybrid CNN--Transformer (one architecture).}
\textbf{TransUNet} \citep{chen2021transunet}: ResNetV2-BiT encoder
\citep{kolesnikov2020} feeding a ViT-B/16 bottleneck, with a cascaded
upsampler decoder using skip connections from encoder stages.

\paragraph{Pure Transformer (one architecture).}
\textbf{Swin-UNet} \citep{cao2022}: Swin-Tiny encoder
\citep{liu2021swin} pretrained on ImageNet-1k, symmetric Transformer
decoder, patch-expanding upsampling.

A complete summary of architectural choices is given in
Table~\ref{tab:archs}. We deliberately exclude all SSM-based variants
(Mamba, Mamba-2, VMamba) from the present benchmark; those are the
subject of a companion study (\PLACE{cite Paper 2 once arxived}).

%==============================================================================
% Paper 1 / T0 — Architecture summary
%==============================================================================
\begin{table*}[t]
\centering
\caption{The nine architectures benchmarked in this paper.
``Pretrained'' indicates whether the encoder is initialised from
ImageNet/BiT pretrained weights or trained from scratch.}
\label{tab:archs}
\small
\setlength{\tabcolsep}{4pt}
\begin{tabular}{llllll}
\toprule
Paradigm           & Architecture       & Encoder              & Decoder           & Pretrained & Params (M) \\
\midrule
Pure CNN           & UNet-V1            & scratch / conv       & transposed conv   & ---        &  31.04 \\
Pure CNN           & UNet-V2            & scratch / conv + SE  & attn-gated skips  & ---        &  33.00 \\
Pure CNN           & UNet-ResNet        & ResNet-34            & U-Net decoder     & ImageNet-1k&  29.07 \\
Pure CNN           & DeepLabV3+         & ResNet-50            & ASPP + light dec. & ImageNet-1k&  40.34 \\
Pure CNN           & nnU-Net            & auto-configured      & deep-supervision  & ---        &  15.62 \\
Pure CNN           & DenseContextU-Net  & dense + context      & dense decoder     & ---        &   2.18 \\
Pure CNN           & FPN-UNet           & ResNet-50            & FPN + U-Net dec.  & ImageNet-1k&  32.16 \\
Hybrid             & TransUNet          & ResNetV2 + ViT-B/16  & cascaded upsmpler & BiT + IN-21k& 102.09 \\
Pure Transformer   & Swin-UNet          & Swin-Tiny            & Swin decoder      & IN-1k      &  41.84 \\
\bottomrule
\end{tabular}
\end{table*}


\subsection{Training protocol}
\label{sec:methods-training}

All nine architectures share a single training protocol:
\begin{itemize}[leftmargin=*,nosep]
\item AdamW optimiser ($\beta_1{=}0.9$, $\beta_2{=}0.999$,
      weight decay $10^{-4}$), initial learning rate $10^{-4}$, cosine
      annealing with warm restarts ($T_0{=}50$ epochs).
\item Composite loss: $\mathcal{L} = \mathcal{L}_{\text{CE}}
      + \mathcal{L}_{\text{Dice}} + 0.1\, \mathcal{L}_{\text{boundary}}$,
      where the boundary term is computed on the signed distance
      transform \citep{kervadec2019}.
\item Image size $256\times 256$ for all CNN and hybrid models;
      $224\times 224$ for Swin-UNet, which requires input divisible by
      $\text{window\_size}\times 2^{n_\text{stages}} = 7\times 32 = 224$.
\item Augmentation: horizontal flip, $\pm 10\deg$ rotation, small elastic
      deformation, mild brightness/contrast jitter.
\item Up to 100 epochs, automatic mixed precision (AMP, FP16), gradient
      clipping at $\|\nabla\|_2 \le 1$ for pretrained-encoder models
      (Swin-UNet, TransUNet) to prevent the early-epoch divergence we
      otherwise observed.
\item Early stopping with patience $20$ on validation mean Dice.
\item Hardware: Google Colab Pro NVIDIA L4 (24~GB) and A100 (40~GB)
      instances.
\end{itemize}

Per-model batch sizes are reported in Table~\ref{tab:archs} and were
selected to maximise GPU utilisation within memory limits. We did
\emph{not} perform per-architecture hyperparameter search: the goal is
explicit cross-architecture comparison under one shared protocol.

\subsection{Evaluation metrics}
\label{sec:methods-metrics}

We report a multi-dimensional set of metrics:

\paragraph{Geometric.}
Per-class Dice and Intersection-over-Union (IoU) for LV-endo, LV-epi,
and LA, computed on the full test set and separately for ED and ES
frames.

\paragraph{Boundary.}
95\textsuperscript{th}-percentile Hausdorff distance (HD95) and average
symmetric surface distance (ASSD), both in \emph{millimetres} using the
per-sample pixel spacing recorded in NIfTI headers. This is consistent
with the official CAMUS challenge protocol and is essential for
cross-paper comparison.

\paragraph{Clinical.}
Biplane Simpson's ejection fraction (EF) computed from the LV-endo
masks at ED and ES across paired 2CH and 4CH views. We report
EF MAE (\%), Pearson correlation $r$ with the ground-truth EF, and
Bland--Altman bias and 95\% limits of agreement \citep{bland1986}.

\paragraph{Phase-stratified.}
All metrics are additionally reported separately for ED and ES frames,
which is informative because ED frames tend to be larger and have clearer
endocardial boundaries.

\paragraph{Quality-stratified robustness.}
We report mean Dice within each CAMUS image-quality grade (Good /
Medium / Poor) \PLACE{\TODO{generate per-quality-grade subsets from
patient metadata}}.

\paragraph{Statistical significance.}
For every pair of models we run a Wilcoxon signed-rank test on per-sample
mean Dice across the test set, with Bonferroni correction across the
$\binom{9}{2} = 36$ pairs. We additionally produce a critical-difference
diagram \citep{demsar2006} illustrating the band of statistically
equivalent top models.

\paragraph{Efficiency.}
Trainable parameters, GFLOPs at $256\times 256$ input (measured with
\texttt{fvcore} on CUDA), single-image inference latency
(mean over $100$ forward passes after $10$ warmup, batch size $1$, FP16),
and peak GPU memory during forward+backward.

\subsection{Reproducibility}
\label{sec:methods-repro}

A single shared evaluation pipeline is used for all models; the same code
path generates the JSON outputs, CSV tables, LaTeX tables, and the
figures in this paper. The training pipeline fixes the random seed
($42$) across PyTorch, NumPy, and Python. Per-model checkpoints, per-test-
patient prediction NIfTI files, and per-sample Dice arrays are released
to allow independent re-analysis without re-running training.
